\section{Conclusion}
\label{sec:conclusion}

Holding the model, the data and the training procedure fixed, we
varied only the reference forecast and the treatment of night hours.
The reference forecast changes both the magnitude and the shape of
reported skill: it accounts for 70 percent of apparent skill at a
six-hour horizon and inverts the trend from monotonically increasing
to a peak at three hours, a shape inversion that holds on a held-out
test year and is backed by significant Diebold-Mariano tests in every
array-horizon cell on both splits. Including night hours deflates
normalised error by approximately 34 percent, closed-form in the
daylight fraction, while leaving skill nearly unchanged. Both effects
exceed the largest difference we measure across five architectures
under an identical protocol, 0.024 in skill score.

\subsection{Recommendations}
\label{sec:conclusion-recommendations}

Three practices, none of them novel, would make results in this
literature comparable. Report a skill score against a stated reference
forecast: against smart persistence alone our results appear to
improve steadily with horizon, against the convex combination of
climatology and persistence recommended by Yang et
al.~\cite{yang2020verification} they do not, and the two references
support opposite conclusions from identical forecasts. Report the
evaluation sample: whether night hours are included, what daylight
threshold is applied, and how many samples each cell contains are each
sufficient to move a reported error by more than the improvement
typically claimed. Report run-to-run variance: none of the 27 papers
we surveyed does, and a single run cannot separate a genuine effect
from retraining noise.

The harness, the 900 run records, and the coded survey with its
supporting evidence are available at an anonymised repository, link
withheld for review.
